\section{Introduction}
\label{sec:intro}

Many people manage money close to the edge: rent clears a few days before
the paycheck, and an overdraft costs more than the purchase that caused it.
A useful forecast of the bank balance is therefore not a single line but a
distribution, and its uncertainty has to be honest. A cone that is too
narrow invites the very overdraft it should warn about; a cone that is too
wide says nothing.

This paper evaluates such a forecaster under two constraints. The first is
\emph{local-first}: a user's transactions stay on their own devices, so the
model is fitted to one account at a time and nothing is pooled across users.
That rules out the strongest generic approach, a global model trained on
many accounts, which we nevertheless include as a reference point. The
second is \emph{reproducibility}: the same feed and seed should give the
same forecast, byte for byte, on each target the application builds.

The forecaster is a three-state Kalman filter \citep{kalman1960} over the
balance, the daily discretionary spending rate and the expected paycheck,
with known paychecks and bills entering as control inputs. Its noise
variances adapt in the manner of covariance matching
\citep{mehra1970,myers1976}. Outlying spending is clipped by a gate and
resampled in the forecast, and a CUSUM detector \citep{page1954} widens the
state covariance after a regime change. A Monte Carlo simulation from the
filtered state gives the cone (\S\ref{sec:method}).

The evaluation followed a written analysis plan (\S\ref{sec:protocol}).
The plan fixed the claims, the pass/fail thresholds, the baselines and the
metrics. Its commit precedes every real-data result, but the earliest
independent timestamp of it does not (\S\ref{sec:timeline}). When the
original model failed on the development split, a revision rule written
before any revision code chose among four candidate variants. The chosen
variant was frozen, and GitHub received the freeze before the main
held-out evaluation. That evaluation ran once, by a command that refuses a
second run without a logged reason. A separately logged run followed, for
ablations the first run had omitted, after the main results were known.

\paragraph{Contributions, including what failed.}
\begin{enumerate}
\item \textbf{A documented, auditable real-data evaluation.} It covers
  \R{ho/data/accounts_total} Berka accounts, split
  \R{c/dev-percent}/\R{c/heldout-percent} by a hash of the account id. It
  reports the commit timeline with GitHub's push times
  (Table~\ref{tab:timeline}), every deviation and post-hoc analysis, and
  a test that forecasts are unchanged when the input is truncated at the
  forecast origin.
\item \textbf{A negative result on CRPS.} On held-out accounts the selected
  model, V2, has lower CRPS than the probabilistic schedule-aware
  baseline. It has higher CRPS at \hB\ and \hC\ days than two baselines: an
  empirical-quantile forecast from the account's own history, and a
  gradient-boosted model pooled across accounts (\S\ref{sec:results}). The
  planned CRPS claim therefore fails against both, under either way of
  averaging the scores.
\item \textbf{A partial result on calibration.} V2's \lC\% intervals meet
  the planned band at every horizon. Its \lD\% interval does at \hB\ and
  \hC\ days and is inconclusive at \hA. Its \lA\% intervals cover far too
  much, so the planned calibration claim fails. V2 misses below its \lC\%
  interval less often than nominal. A conformal recalibration of the
  original model (V1) is the best-calibrated variant but lost the
  selection on CRPS.
\item \textbf{Regime detection measured two ways.} On real data, the
  comparison with the detector removed is descriptive only, because its
  ``after a change'' stratum is defined by the detector's own alarms. A
  synthetic sweep with known change dates shows a large benefit after large
  changes on model-matched data. On the realistic generator the effect on
  CRPS is imprecise, and the detector costs CRPS when a downward change
  goes undetected (\S\ref{sec:regime}).
\item \textbf{Bit identity, falsified and then fixed.} The frozen engine
  was not byte-identical across native code and WebAssembly. We give the
  cause, a post-freeze fix and the tested matrix (\S\ref{sec:identity}).
\end{enumerate}

The method is a composition of known parts, and we do not claim novelty
for the parts. The question the paper answers is narrower. On sparse,
booked bank accounts, does careful state-space modelling give lower CRPS
than simple alternatives built from the same account's history? On Berka
it does against a schedule-aware baseline, and it does not against an
empirical distribution of the account's own balance changes.
